% Options for packages loaded elsewhere
\PassOptionsToPackage{unicode}{hyperref}
\PassOptionsToPackage{hyphens}{url}
\PassOptionsToPackage{dvipsnames,svgnames,x11names}{xcolor}
\documentclass[
  10pt,
  fontset=fandol]{ctexart}
\usepackage{xcolor}
\usepackage[margin=16mm]{geometry}
\usepackage{amsmath,amssymb}
\setcounter{secnumdepth}{-\maxdimen} % remove section numbering
\usepackage{iftex}
\ifPDFTeX
  \usepackage[T1]{fontenc}
  \usepackage[utf8]{inputenc}
  \usepackage{textcomp} % provide euro and other symbols
\else % if luatex or xetex
  \usepackage{unicode-math} % this also loads fontspec
  \defaultfontfeatures{Scale=MatchLowercase}
  \defaultfontfeatures[\rmfamily]{Ligatures=TeX,Scale=1}
\fi
\usepackage{lmodern}
\ifPDFTeX\else
  % xetex/luatex font selection
\fi
% Use upquote if available, for straight quotes in verbatim environments
\IfFileExists{upquote.sty}{\usepackage{upquote}}{}
\IfFileExists{microtype.sty}{% use microtype if available
  \usepackage[]{microtype}
  \UseMicrotypeSet[protrusion]{basicmath} % disable protrusion for tt fonts
}{}
\makeatletter
\@ifundefined{KOMAClassName}{% if non-KOMA class
  \IfFileExists{parskip.sty}{%
    \usepackage{parskip}
  }{% else
    \setlength{\parindent}{0pt}
    \setlength{\parskip}{6pt plus 2pt minus 1pt}}
}{% if KOMA class
  \KOMAoptions{parskip=half}}
\makeatother
\usepackage{longtable,booktabs,array}
\newcounter{none} % for unnumbered tables
\usepackage{calc} % for calculating minipage widths
% Correct order of tables after \paragraph or \subparagraph
\usepackage{etoolbox}
\makeatletter
\patchcmd\longtable{\par}{\if@noskipsec\mbox{}\fi\par}{}{}
\makeatother
% Allow footnotes in longtable head/foot
\IfFileExists{footnotehyper.sty}{\usepackage{footnotehyper}}{\usepackage{footnote}}
\makesavenoteenv{longtable}
\setlength{\emergencystretch}{3em} % prevent overfull lines
\providecommand{\tightlist}{%
  \setlength{\itemsep}{0pt}\setlength{\parskip}{0pt}}
\usepackage{amsmath,amssymb,mathtools,needspace}
\xeCJKDeclareCharClass{CJK}{"2032,"0400->"04FF,"2160->"216B,"2460->"2473,"25A0,"25C7,"25CB,"25CF}
\IfFontExistsTF{Arial Unicode MS}{\xeCJKsetup{AutoFallBack=true}\setCJKfallbackfamilyfont{\CJKrmdefault}{Arial Unicode MS}}{}
\setcounter{secnumdepth}{0}
\setlength{\emergencystretch}{2em}
\allowdisplaybreaks
\usepackage{bookmark}
\IfFileExists{xurl.sty}{\usepackage{xurl}}{} % add URL line breaks if available
\urlstyle{same}
\hypersetup{
  pdftitle={偏差比中传出好``消息''},
  colorlinks=true,
  linkcolor={Maroon},
  filecolor={Maroon},
  citecolor={Blue},
  urlcolor={Blue},
  pdfcreator={LaTeX via pandoc}}

\title{偏差比中传出好``消息''}
\author{}
\date{}

\begin{document}
\maketitle

\subsubsection{12.3.2　偏差比中传出好``消息''}\label{ux504fux5deeux6bd4ux4e2dux4f20ux51faux597dux6d88ux606f}

再换一种视角考察前述数据校正技术。针对无穷逼近过程 \(\{S_n\}\)，能否设计更``好''的逼近过程 \(\{\widehat S_n\}\)，使之具有更快的收敛速度呢？这就是逼近加速问题。

欲使\textbf{校正公式}（12.3.1）成为逼近加速公式，校正因子 \(\omega\) 应当是驾驭逼近过程的某个数学不变量，即与逼近过程相关的某个普适常数，称之为\textbf{加速因子}。前已指出，这个常数 \(\omega\) 应当在 0 与 1 之间，即 \(0<\omega<1\)。

数学常数通常采取比率的形式，刘徽指出，比率的本意是相关事物的比例关系。

在二分割圆过程中，需要考虑什么样的``相关量''呢？

《割圆术》称偏差 \(\Delta_n=S_{2n}-S_n\) 为\textbf{差幂}，刘徽特别关注\textbf{幂率} \(\delta_n=\Delta_n/\Delta_{2n}\)。在割圆计算过程中，利用差幂 \(\Delta_n\) 计算幂率 \(\delta_n\)，割圆术实际上已造出下列数据表 12.1。

\textbf{表 12.1}

\begingroup
\renewcommand{\arraystretch}{2.2}

{\def\LTcaptype{none} % do not increment counter
\begin{longtable}[]{@{}rrrr@{}}
\toprule\noalign{}
\(n\) & \(S_n\) & \(\Delta_n\) & \(\delta_n\) \\
\midrule\noalign{}
\endhead
\bottomrule\noalign{}
\endlastfoot
12 & 300 & \(10\dfrac{364}{625}\) & 3.95 \\
24 & \(310\dfrac{364}{625}\) & \(2\dfrac{425}{625}\) & 3.99 \\
48 & \(313\dfrac{164}{625}\) & \(\dfrac{420}{625}\) & 4.00 \\
96 & \(313\dfrac{584}{625}\) & \(\dfrac{105}{625}\) & \\
192 & \(314\dfrac{64}{625}\) & & \\
\end{longtable}
}

\endgroup

从这些数据中能获得什么样的信息呢？无需用高深的数学知识进行理论分析，\textbf{直接观察幂率 \(\delta_n\) 的数据表（见表 12.1）即可发现，幂率几乎为定值 4。}

\textbf{刘徽由此发现了一个奇妙的事实：表面上杂乱无章的数据 \(S_n\) 中竟蕴藏着极其鲜明的规律性。}

\end{document}
